\begin{resumo}[Abstract]
The occurrence of adverse events (AEs) in hospital environments represents one of the most critical global challenges in public healthcare, directly impacting patient morbidity and mortality and generating significant socioeconomic burdens. Traditionally, healthcare institutions rely on passive, voluntary reporting systems, which capture less than 10\% of actual harm. In response to this limitation, the Institute for Healthcare Improvement (IHI) developed the Global Trigger Tool (GTT) methodology, structured around the active, retrospective review of electronic medical records guided by clinical ``triggers''. This work presents the proposal for SIGEA-GTT (Intelligent System for Adverse Event Management – Global Trigger Tool), an interdisciplinary computational platform designed through a collaborative effort between the Department of Computing (DCOMP) and the Department of Nursing at the Federal University of Sergipe (UFS). The system integrates the IHI-GTT methodology (6 specialized modules, 53 standardized triggers, and NCC MERP harm severity classification from A to I) with established quality engineering tools (Ishikawa Diagram, GUT Matrix, 5W2H Action Plan, and PDCA/PDSA Cycle). The architectural framework is structured around modular backend services in Java 21 LTS and Spring Boot 3.3, relational persistence with JSONB in PostgreSQL 16, and an ergonomic frontend developed in Angular 18 with Standalone Components, Signals, and TailwindCSS, tailored for high data-density clinical workstations. The project also encompasses an innovative educational simulation engine for academic and clinical training at the University Hospital of UFS (HU-UFS), governed by a strict Quality Gate enforcing 100\% automated test coverage.

\textbf{Keywords}: Patient Safety. Global Trigger Tool. Adverse Events. Quality Management. Healthcare Software Engineering. Federal University of Sergipe.
\end{resumo}
